\chapter{Codice Utilizzato}
\label{app:B}

Questa appendice raccoglie i listati richiamati nei capitoli principali. Nel
testo sono mantenute le scelte progettuali, le decisioni implementative e il
razionale; i frammenti di codice sono spostati qui per non interrompere la
lettura.

\section{Listati richiamati nel Capitolo 2}

\begin{lstlisting}[language=Python, frame=single, breaklines=true, caption={Struttura della pipeline di ingestione Haystack 2.x (pseudocodice).}, label={lst:ingest_pipeline}]
from haystack import Pipeline
from haystack.components.writers import DocumentWriter
from haystack_integrations.components.embedders.fastembed import (
    FastembedDocumentEmbedder,
    FastembedSparseDocumentEmbedder,
)
from haystack_integrations.document_stores.qdrant import QdrantDocumentStore

qdrant_store = QdrantDocumentStore(url="http://qdrant:6333",
                                   index="i3guild_knowledge",
                                   embedding_dim=768,
                                   use_sparse_embeddings=True)

indexing_pipeline = Pipeline()
indexing_pipeline.add_component("dense_embedder", FastembedDocumentEmbedder(
    model="sentence-transformers/paraphrase-multilingual-mpnet-base-v2"))
indexing_pipeline.add_component("sparse_embedder", FastembedSparseDocumentEmbedder(
    model="Qdrant/bm25"))
indexing_pipeline.add_component("writer", DocumentWriter(
    document_store=qdrant_store,
    policy=DuplicatePolicy.OVERWRITE))

indexing_pipeline.connect("dense_embedder.documents", "sparse_embedder.documents")
indexing_pipeline.connect("sparse_embedder.documents", "writer.documents")

indexing_pipeline.run({"dense_embedder": {"documents": raw_docs}})
\end{lstlisting}

\begin{lstlisting}[frame=single, breaklines=true, caption={Esempio di request body per l'endpoint \texttt{/rag/retrieve}.}, label={lst:api_request}]
{
  "query":           "vorrei fondare una gilda sulla musica",
  "top_k":           20,
  "score_threshold": 0.5,
  "max_documents":   10,
  "gap_threshold":   0.08,
  "filters":         null
}
\end{lstlisting}

\begin{lstlisting}[frame=single, breaklines=true, caption={Struttura della response dell'endpoint \texttt{/rag/retrieve}.}, label={lst:api_response}]
{
  "documents":        [ { "content": "...", "score": 0.85,
                          "meta": { "guild_name": "...",
                                    "participants": [...] } } ],
  "total_candidates": 15,
  "threshold_used":   0.5,
  "cutoff_reason":    "relative_gap",
  "topic_query":      "musica"
}
\end{lstlisting}

\section{Listati richiamati nel Capitolo 3: pipeline di ingestione}

\begin{lstlisting}[language=SQL, frame=single, breaklines=true, caption={Campi principali estratti da PostgreSQL per l'ingestione.}, label={lst:postgres_extract_fields}]
SELECT
    g.id AS guild_id,
    g.name AS guild_name,
    g.summary,
    g.goals,
    g.opportunities,
    g.outcome,
    g.risks,
    g."otherResources" AS other_resources,
    g.budget,
    g."achievedResultsDescription" AS achieved_results_description,
    g.journal,
    g."isIndividualGuild" AS is_individual,
    g."participationMode" AS participation_mode,
    e.id AS epoch_id,
    e."startDate" AS epoch_start,
    e."endDate" AS epoch_end
FROM "Guild" g
LEFT JOIN "Epoch" e ON g."epochId" = e.id
\end{lstlisting}

\begin{lstlisting}[language=Python, frame=single, breaklines=true, caption={Schema semplificato della costruzione del documento testuale.}, label={lst:create_document_simplified}]
parts = [f"[GUILD: {guild_name}]"]

if row.get("summary"):
    parts.append(f"[SOMMARIO] {clean_html(row.get('summary'))}")
if row.get("goals"):
    parts.append(f"[OBIETTIVI] {clean_html(row.get('goals'))}")
if row.get("risks"):
    parts.append(f"[RISCHI] {clean_html(row.get('risks'))}")
if participants:
    parts.append(f"[PARTECIPANTI] {', '.join(participants)}")

full_text = "\n".join(parts)
\end{lstlisting}

\begin{lstlisting}[language=Python, frame=single, breaklines=true, caption={Pipeline Haystack usata per embedding e scrittura.}, label={lst:embedding_pipeline_impl}]
document_store = QdrantDocumentStore(
    url=QDRANT_URL,
    index=QDRANT_COLLECTION_NAME,
    embedding_dim=EMBEDDING_DIMENSIONS,
    similarity="cosine",
    recreate_index=RECREATE_INDEX,
    use_sparse_embeddings=True,
)

dense_embedder = FastembedDocumentEmbedder(
    model=FASTEMBED_DENSE_MODEL,
    cache_dir=FASTEMBED_CACHE_DIR,
)
sparse_embedder = FastembedSparseDocumentEmbedder(
    model=FASTEMBED_SPARSE_MODEL,
    cache_dir=FASTEMBED_CACHE_DIR,
)

writer = DocumentWriter(document_store=document_store, policy=policy)
indexing_pipeline = Pipeline()
indexing_pipeline.add_component("dense_embedder", dense_embedder)
indexing_pipeline.add_component("sparse_embedder", sparse_embedder)
indexing_pipeline.add_component("writer", writer)
indexing_pipeline.connect("dense_embedder.documents", "sparse_embedder.documents")
indexing_pipeline.connect("sparse_embedder.documents", "writer.documents")
\end{lstlisting}

\section{Listati richiamati nel Capitolo 3: retrieval e generazione}

\begin{lstlisting}[language=Python, frame=single, breaklines=true, caption={Schema Pydantic della richiesta di retrieval.}, label={lst:retrieve_request}]
class RetrieveRequest(BaseModel):
    query: str
    top_k: int = 20
    score_threshold: float = 0.7
    max_documents: int = 10
    gap_threshold: float = 0.08
    filters: dict | None = None
\end{lstlisting}

\begin{lstlisting}[language=Python, frame=single, breaklines=true, caption={Pipeline Haystack per la ricerca ibrida.}, label={lst:query_pipeline_impl}]
document_store = QdrantDocumentStore(
    url=QDRANT_URL,
    index=QDRANT_COLLECTION_NAME,
    embedding_dim=768,
    similarity="cosine",
    recreate_index=False,
    use_sparse_embeddings=True,
)
text_embedder = FastembedTextEmbedder(
    model=FASTEMBED_DENSE_MODEL,
)
sparse_text_embedder = FastembedSparseTextEmbedder(
    model=FASTEMBED_SPARSE_MODEL,
)
retriever = QdrantHybridRetriever(document_store=document_store)

pipeline = Pipeline()
pipeline.add_component("text_embedder", text_embedder)
pipeline.add_component("sparse_text_embedder", sparse_text_embedder)
pipeline.add_component("retriever", retriever)
pipeline.connect("text_embedder.embedding", "retriever.query_embedding")
pipeline.connect("sparse_text_embedder.sparse_embedding",
                 "retriever.query_sparse_embedding")
\end{lstlisting}

\begin{lstlisting}[language=Python, frame=single, breaklines=true, caption={Logica semplificata di taglio dinamico dei risultati.}, label={lst:dynamic_filtering}]
scored_docs.sort(key=lambda x: x[0], reverse=True)
above_threshold = [
    d for score_value, d in scored_docs
    if score_value >= score_threshold
]

filtered = [above_threshold[0]]
for i in range(1, len(above_threshold)):
    prev = float(getattr(above_threshold[i - 1], "score", 0.0) or 0.0)
    cur = float(getattr(above_threshold[i], "score", 0.0) or 0.0)
    relative_drop = (prev - cur) / prev if prev > 0 else 1.0
    if relative_drop > gap_threshold:
        cutoff_reason = "relative_gap"
        break
    filtered.append(above_threshold[i])
\end{lstlisting}

\begin{lstlisting}[language=Python, frame=single, breaklines=true, caption={Formato del messaggio utente aumentato con contesto RAG.}, label={lst:augmented_prompt}]
final_message = (
    f"{message}\n\n"
    "--- CONTESTO RAG ---\n"
    "Usa queste informazioni per rispondere.\n"
    f"{final_context}\n"
    "--- FINE CONTESTO ---"
)
\end{lstlisting}

\begin{lstlisting}[frame=single, breaklines=true, caption={Esempio di eventi NDJSON prodotti dalla chat standard.}, label={lst:ndjson_standard}]
{"type": "rag", "data": {"documents": [...], "cutoff_reason": "relative_gap"}}
{"type": "stream_start", "model": "qwen2.5:1.5b"}
{"type": "token", "content": "La"}
{"type": "token", "content": " gilda"}
{"type": "stream_end"}
\end{lstlisting}

\section{Listati richiamati nel Capitolo 3: frontend e generazione bozza}

\begin{lstlisting}[frame=single, breaklines=true, caption={Esempio di eventi NDJSON prodotti dalla generazione bozza.}, label={lst:draft_ndjson}]
{"type": "status", "message": "generation_started", "promptLength": 4200}
{"type": "status", "message": "generation_in_progress"}
{"type": "draft", "source": "llm", "draft": {"name": "..."}}
\end{lstlisting}

\begin{lstlisting}[frame=single, breaklines=true, caption={Chiamata TypeScript per generazione bozza senza retrieval.}, label={lst:draft_generation_call}]
const ragResponse = await generateRAGChatReply(prompt, [], {
    model,
    useRag: false,
    jsonMode: true,
    options: { num_predict: -1, timeout_seconds: 0 },
});
\end{lstlisting}

\begin{lstlisting}[frame=single, breaklines=true, caption={Chiamata TypeScript al retrieval RAG.}, label={lst:typescript_search_rag}]
const response = await fetch(`${RAG_API_URL}/rag/retrieve`, {
    method: "POST",
    headers: { "Content-Type": "application/json" },
    body: JSON.stringify({
        query,
        filters,
        top_k: 20,
        score_threshold: options?.scoreThreshold ?? 0.7,
        max_documents: options?.maxDocuments ?? 10,
    }),
});
\end{lstlisting}
